% Source-aligned English for printed p.129.
\SourcePage{129}

\LayerTarget{tohoku:section:1.5:p1}
\noindent in $C$, is \textit{right exact}. It is even exact if $I$ is finite,
but not necessarily if $I$ is infinite, since the direct sum of an infinite
family of monomorphisms need not be a monomorphism, as we observed in No.~1
(for the dual situation). Hence the following axiom:

\LayerTarget{tohoku:section:1.5:AB4}
\noindent\textbf{AB 4)} \textit{Axiom AB 3) holds, and the direct sum of a
family of monomorphisms is a monomorphism.}

The following axiom is strictly stronger than AB 4):

\LayerTarget{tohoku:section:1.5:AB5}
\noindent\textbf{AB 5)} \textit{Axiom AB 3) holds, and if $(A_i)_{i\in I}$ is
an increasing directed family of subobjects of an object $A\in C$, and $B$ is
any subobject of $A$, then}

\[
  \left(\sum_i A_i\right)\cap B=\sum_i(A_i\cap B).
\]

\noindent (In accordance with usage, $\sum A_i$ denotes the supremum of the
$A_i$, and $P\cap Q$ the infimum of the subobjects $P$ and $Q$ of $A$.)

\LayerTarget{tohoku:section:1.5:AB5:equivalent}
AB 5) can also be expressed as follows: \textit{AB 3) holds, and if $A\in C$
is the supremum of an increasing directed family of subobjects $A_i$, and if
for every $i$ one is given a morphism $u_i:A_i\longrightarrow B$ such that,
whenever $A_i\subset A_j$, $u_i$ is induced by $u_j$, then there exists a
morphism $u$} (evidently unique) \textit{from $A$ to $B$ that induces the
$u_i$}. Finally, note the following axiom, which strengthens AB 5) still
further, but which we shall not need to use in this work:

\LayerTarget{tohoku:section:1.5:AB6}
\noindent\textbf{AB 6)} \textit{Axiom AB 3) holds, and for every $A\in C$ and
every family $(B^j)_{j\in J}$ of increasing directed families
$B^j=(B_i^j)_{i\in I_j}$ of subobjects $B_i^j$ of $A$, one has}

\[
 \bigcap_{j\in J}\left(\sum_{i\in I_j}B_i^j\right)
 =\sum_{(i_j)\in\prod_{j\in J}I_j}
   \left(\bigcap_{j\in J}B_{i_j}^j\right).
\]

\noindent (This axiom implicitly includes the existence of the infimum of an
arbitrary family of subobjects of $A$.)

\LayerTarget{tohoku:section:1.5:p2}
We leave to the reader the task of stating the dual axioms AB 3*), AB 4*),
AB 5*), and AB 6*), concerning infinite products. As an example, note that the
category of abelian groups (or, more generally, the category of modules over a
fixed unital ring) satisfies, for direct sums, the strongest axiom AB 6); it
also satisfies AB 3*) and AB 4*), but not AB 5*). The facts are reversed for
the dual category, which by Pontryagin duality is isomorphic to the category of
compact topological abelian groups. (This shows that AB 5*) is not a
consequence of AB 4*), and hence that AB 5) is likewise not a consequence of
AB 4).) The abelian category of sheaves of abelian groups on a fixed
topological space $X$ satisfies AB 5) and AB 3*), but not AB 4*), since we have
already observed that a product of surjective morphisms need not be a
surjective morphism. Finally, note that if $C$ is a category satisfying both
AB 5) and AB 5*), then $C$ is reduced to null objects (for one then sees easily
that, for $A\in C$, the canonical morphism
$A^{(I)}\longrightarrow A^I$ is an isomorphism, and one verifies that this is
possible only if $A$ is null).

\LayerTarget{tohoku:section:1.5:p3}
The preceding axioms will be especially useful to us in studying inductive and
projective limits, which we shall need in order to give conditions

\TohokuContinuation{130}
